\subsection{Entity-Based Graph Construction}
Entity-based graphs leverage named entities (e.g., individuals, organizations, locations) to encode relationships between textual units, enabling applications like topic tracking, event detection, and cross-platform analysis. Syed et al. \citep{syed2021unified} proposed a tripartite graph structure that combines tweets, news articles, and entities, extracting semantic relations between heterogeneous sources. Their approach used entity co-occurrence to link Twitter tweets and web news, with superior event detection performance compared to traditional keyword-based methods.
Spitz and Gertz \citep{spitz2018entity} developed an entity-based method for news stream aggregation, constructing graphs with entities as vertices and edges as co-occurrence in news. Their method enabled the detection of evolving news events by monitoring relationships between entities over time. Manaskasemsak et al. \citep{manaskasemsak2024entity} built co-occurrence entity graphs from tweets to group tweets into events with high accuracy in detecting localized events like protests or natural disasters.
Recent work further expands the utility of entity graphs for social media analysis. Bodaghi et al. (2025) analyzed Twitter news entity networks to uncover macro-scale structural features in global news dissemination, underscoring the value of entity graph structures for understanding information flow in evolving news ecosystems \citep{bodaghi2025decoding}. Additionally, emerging approaches for security and event detection on Twitter integrate entity patterns into embedding and detection frameworks to improve precision and coverage of security-related events \citep{cui2024enhancing}. These developments affirm the robustness of entity-based graph representations and extend their applicability to dynamic and temporally sensitive social media scenarios.

\subsection{Social Network Analysis on Twitter}
Social Network Analysis (SNA) of Twitter provides the structural properties of communication networks, including topology, influence dynamics, and information spread. Himelboim et al. \citep{himelboim2017classifying} performed SNA of Twitter conversation networks through measures like density and modularity, categorizing them as hub-and-spoke (broadcast) and clustered (community) structures. They found that hub-and-spoke networks, dominated by influential users or news outlets, dominate breaking news, while clustered networks emerge in polarized debates. Logan et al. \citep{logan2023modeling} generalized SNA to multilayer networks, using PageRank-based influencer detection to reveal hidden patterns across interaction layers. Kwak et al. \citep{kwak2010what} also conducted pioneering SNA analysis of Twitter, analyzing follower-followee networks and information spreading.
Recent advances have increasingly incorporated graph neural networks (GNNs) into Twitter SNA to model complex relational patterns and improve detection tasks. For example, Patel et al. (2024) used GNNs with data augmentation for rumour detection in Twitter datasets, showing large performance improvements over baseline classifiers \citep{patel2024rumour}. Advanced frameworks like ALETHEIA (2025) employ GNNs with temporal link prediction to detect malicious influence campaigns by modeling troll and user interactions at scale \citep{saeed2025aletheia}. These studies highlight the growing trend of combining topological information with deep learning to uncover influence dynamics and network roles.
In addition to GNN-based work, the SNA literature continues to innovate on influence and engagement modeling. New research on influential user identification explores temporal interaction patterns and indirect influence pathways to improve the identification of key nodes in information spread \citep{hassanpour2026identification}. These approaches complement traditional centrality measures (e.g., degree, betweenness), providing richer tools for structural analysis and engagement metrics in Twitter interaction graphs. Taken together, these studies illustrate SNA’s capacity to model Twitter’s flexibility and dynamism, offering methodological insights for degree and centrality analysis in networked conversation systems.
